% !TeX program = pdflatex
\documentclass[11pt,a4paper]{article}
\usepackage[margin=27mm]{geometry}
\usepackage{amsmath,amssymb,amsthm}
\usepackage[T1]{fontenc}
\usepackage{lmodern}
\usepackage[hidelinks]{hyperref}
\hypersetup{pdftitle={Arithmetic obstructions and finite tests for the Li--Liu spectral hypotheses},pdfauthor={Codex}}
\newtheorem{theorem}{Theorem}[section]
\newtheorem{lemma}[theorem]{Lemma}
\newtheorem{corollary}[theorem]{Corollary}
\newtheorem{proposition}[theorem]{Proposition}
\newcommand{\SSg}{\mathfrak S}
\newcommand{\one}{\mathbf 1}
\newcommand{\dd}{\,\mathrm d}
\DeclareMathOperator{\Rea}{Re}
\DeclareMathOperator{\supp}{supp}
\allowdisplaybreaks[1]
\setlength{\emergencystretch}{1em}
\title{Arithmetic obstructions and finite tests for the\\ Li--Liu spectral hypotheses}
\author{Codex}
\date{27 November 2026}
\begin{document}
\maketitle
\begin{abstract}
Li and Liu deduce a weighted binary Goldbach asymptotic from a prescribed zero-free region for the spectral function of a residual operator. Under Dirichlet GRH we construct an arithmetic sequence on which the zero-free region fails. Any admissible prescription can be modified on an arbitrarily thin subsequence of these records to produce a counterexample. In particular, a failing prescription agrees with $\log N$ off a set containing at most $1+\lfloor\log_2 X\rfloor$ integers up to $X$, and its width is smaller than every fixed inverse power of $\log N$. We give an exact finite-matrix criterion for the unmodified prescription and for the existence of any admissible zero-free region. On the first physical scale, decay of the normalized quadratic form on its entire test family is equivalent to a uniform Hardy--Littlewood asymptotic for ordinary ordered prime pairs with a sharp cutoff in the first summand. An unconditional total variation estimate removes both von Mangoldt weights. We also prove a two-prime asymptotic on both sum and difference congruences at moduli of size comparable to the summation length divided by the cube of its logarithm. The construction also gives an independent derivation of the failure of Hypothesis R under Dirichlet GRH.
\end{abstract}

\section{The hypotheses and the main results}

Let $\Lambda$ be the von Mangoldt function. For even $N\ge4$, write
\[
 C_2=\prod_{p>2}\left(1-\frac{1}{(p-1)^2}\right),\qquad
 \SSg(N)=2C_2\prod_{\substack{p\mid N\\p>2}}\frac{p-1}{p-2},
 \qquad \delta_N^2=\frac{\SSg(N)(\log N)^3}{N}.
\]
Products indexed by $p$ are over primes. Write $G(N)$ for GRH in the open critical strip for primitive Dirichlet characters of conductors dividing $N$. Let $G$ denote the assertion that $G(N)$ holds for every sufficiently large even $N$. Each positive conductor divides arbitrarily large even integers, so $G$ is equivalent to GRH for all primitive Dirichlet characters. The arithmetic obstructions below use $G$. The equivalences concerning individual moduli retain the local condition $G(N)$.

For $N>e^{10}$ satisfying $G(N)$, let $T_N$ be the Li--Liu compact selfadjoint residual operator on the Hardy strip $a_N<\Rea s<b$, where $a_N=1/\log N$ and $b=1/8$. We use their normalization $c=1$ throughout. Its nonzero eigenvalues $\mu_\nu$, counted with multiplicity, define
\[
 X_N(q)=\prod_\nu\left(1-\frac{4\mu_\nu^2}{\delta_N^2}\,q(1-q)\right)
 \qquad(q\in\mathbb H).
\]
Hypothesis R requires all its zeros to have real part $1/2$ for every sufficiently large eligible even $N$. Here and below, \emph{eligible} means that $G(N)$ holds and the strip parameters are valid. A prescription $h$ is admissible if it is positive on even $N\ge4$, tends to infinity along the even integers, and eventually satisfies
\[
  \eta_h(N)=\frac{h(N)^2(\log N)^3}{N\SSg(N)}\le \tfrac14 .
\]
For each fixed $h$, Hypothesis $\mathrm S(h)$ asserts admissibility and, for all sufficiently large eligible even $N$, the exclusion
\begin{equation}\label{eq:free}
 X_N(q)\ne0 \qquad\text{whenever }\Rea q>1-\eta_h(N).
\end{equation}
Under $G$, Li and Liu's Theorem 1.1 deduces a weighted binary Goldbach asymptotic from any one such $h$ \cite{LL}. Thus refuting $\mathrm S(h)$ for a family of functions and refuting $\exists h\,\mathrm S(h)$ are distinct conclusions.

\begin{theorem}[Sparse changes of prescription]\label{th:sparse}
Assume $G$. There exist strictly increasing even integers $r_0,r_1,\dots$ and an admissible $h_*$, defined by the arithmetic tests in Section~\ref{sec:records}, with the following property. Given any admissible $h_0$ and any strictly increasing $p:\mathbb N\to\mathbb N$, put $E_p=\{r_{p(j)}:j\ge0\}$ and
\[
 \widetilde h(N)=
 \begin{cases}h_*(N),&N\in E_p,\\ h_0(N),&N\notin E_p.\end{cases}
\]
Then $\widetilde h$ is admissible and $\mathrm S(\widetilde h)$ fails. If $p(j)\ge2^j$, then
$\#(E_p\cap[1,X])\le 1+\lfloor\log_2 X\rfloor$ for every integer $X\ge1$.
For $h_0(N)=\log N$ and every fixed integer $d\ge0$,
$(\log N)^d\eta_{\widetilde h}(N)\to0$.
\end{theorem}

The records are defined without using a norm, eigenvalue, or zero of $T_N$. The sequence $p$ may grow as fast as desired. The argument gives no conclusion about $\mathrm S(\log)$ from equality of the two prescriptions outside $E_p$.

\begin{theorem}[Arithmetic norm growth]\label{th:growth}
Assume $G$. For each integer $d\ge0$ there are arbitrarily large even $N$ such that $N\|T_N\|^2>(\log N)^d$. Hypothesis R fails. If for some fixed $j\ge0$ an admissible prescription eventually satisfies
$\eta_h(N)\ge 1/(4(\log N)^j)$, then $\mathrm S(h)$ fails. On a single increasing sequence $N_i$ of even integers there are real zeros $x_i<1$ of $X_{N_i}$ such that $(\log N_i)^d(1-x_i)\to0$ for every fixed $d\ge0$.
\end{theorem}

Yilin Liu previously obtained an obstruction to R and to prescribed inverse-logarithmic widths by detecting opposite prime shifts with a Maier matrix \cite{YL}. We use the opposite shifts, derive the required operator pairing, and construct the arithmetic discrepancy by finite inclusion--exclusion and the ordinary prime number theorem. Section~\ref{sec:samples} treats the two remaining assertions directly through the entries of a fixed sequence of sampling matrices. Section~\ref{sec:limit} bounds the lower estimate obtainable from one short-shift test at their scales.

\section{Spectral thresholds and an arithmetic matrix element}\label{sec:entry}

For $0<\eta<1/2$ the norm criterion for an open zero-free half-space is
\begin{equation}\label{eq:dict}
 X_N(q)\ne0\ (\Rea q>1-\eta)
 \quad\Longleftrightarrow\quad
 4\eta(1-\eta)\|T_N\|^2\le\delta_N^2 .
\end{equation}
Indeed a nonzero norm is attained in absolute value by an eigenvalue. On each complex slice the corresponding roots are
$\bigl(1\pm\sqrt{1-\delta_N^2/\mu_\nu^2}\bigr)/2$.
Square summability of the eigenvalues ensures that the product has no other zeros. A strict violation of the right-hand bound gives a real root in $(1-\eta,1)$. Equality puts the extremal root on the boundary. The same argument gives the criterion $\|T_N\|\le\delta_N$ for R at $N$ \cite[Theorem 6.1]{LL}.

We specify the kernel used throughout. Put $L_\chi(s)=(L'/L)(1/2+s,\chi)$, $w=s+z$, and
\begin{align*}
 D_N(w)&=\SSg(N)\zeta(w)\prod_{\substack{p>2\\p\nmid N}}
                    (1+p^{-w}/(p-2)),\\
 \mathcal B(s,z)&=B(1/2-s,w)+B(1/2-z,w)-B(1/2-s,1/2-z),\\
 \mathcal E_N(s,z)&=\frac{2}{\varphi(N)}
       \sum_{\substack{\chi\bmod N\\\chi(-1)=-1}} L_\chi(s)\overline{L_\chi(\bar z)}
       -\sum_{(n,N)=1}\frac{\Lambda(n)^2}{n^{1+w}}-N^{-w}D_N(w)\mathcal B(s,z).
\end{align*}
Here $B$ is Euler's beta function. The weighted kernel is
$K_N(s,z)=\mathcal E_N(s,\bar z)/((1+s)(1+\bar z))$. The Hardy norm uses the sum of the two boundary $L^2(\mathbb R,\dd t)$ integrals. In this normalization $K_N(s,z)=\langle T_Nk_z,k_s\rangle$ for the evaluation vectors $k_s$ and the inner product linear in its first variable.

Let $1\le H\le P<N/20$. A test is a real $b_0\in C^\infty_c(\mathbb R)$ supported in $[P,2P]$ and in an interval of length $H$, with $|b_0|\le1$ and $H|b_0'|\le1$. Write
\[
 B_N(b_0)=\sum_{\substack{m<N\\(m,N)=1}}\Lambda(m)b_0(m)
                 \bigl(\Lambda(N-m)-\Lambda(N+m)\bigr).
\]
\begin{lemma}\label{lm:transfer}
For an absolute constant $C$ and every sufficiently large eligible even $N$,
\begin{equation}\label{eq:transfer}
 H B_N(b_0)^2\le C P^2 N\|T_N\|^2
\end{equation}
for every test.
\end{lemma}
\begin{proof}
Set $M=2N$ and $f(u)=b_0(Mu)$. Choose a fixed smooth $g$, symmetric about $1/2$, taking values in $[0,1]$, supported in $[3/8,5/8]$, and equal to one on $[7/16,9/16]$. For a real smooth $v$ compactly supported in $(0,1)$ set
\[
 w_v(n)=M^{-1/2}\Lambda(n)v(n/M),\qquad
 J_v(y)=2\int_{\mathbb R}v(u)v(u+y)\dd u-\int_{\mathbb R}v(u)v(y-u)\dd u.
\]
For a Hardy function $F$ let $x_v$ represent the linear functional
\[
 F\longmapsto\frac1{2\pi i}\int_{\Rea s=1/9}
 M^s\widehat{\sqrt u\,v(u)}(s)(1+s)F(s)\dd s,\qquad
 \widehat f(s)=\int_0^\infty f(u)u^{s-1}\dd u .
\]
Pairing each of the three kernels with this functional twice gives
\begin{align}
 \langle T_N x_v,x_v\rangle={}&
 \frac{2}{\varphi(N)}\sum_{\chi(-1)=-1}\left|\sum_{n<M}w_v(n)\chi(n)\right|^2
 -\sum_{\substack{n<M\\(n,N)=1}}w_v(n)^2
 -\sum_{\ell\ge1}\SSg(\ell N)J_v(\ell/2).
 \label{eq:pair}
\end{align}
For clarity, in the character term Mellin inversion first uses the absolutely convergent expansion of $-L'/L$. The move to $\Rea s=1/9$ crosses no pole for an odd character under $G(N)$. The diagonal is inverted termwise. For the last term use $w=s+z$ in the two Mellin integrals. The inner beta integral is the Mellin transform of $J_v$. The quotient of the sampling scales gives $2^w$. On $\Rea w>1$, $D_N(w)=\sum_{\ell\ge1}\SSg(\ell N)\ell^{-w}$, so inversion samples at $\ell/2$. The pole at $w=1$ is removable in the pairing because $\int_0^\infty J_v(y)\dd y=0$. This follows by interchanging the two integrals and using symmetry in their two arguments. Smooth compact support gives rapid vertical decay in the contour moves.

Polarize \eqref{eq:pair} using $v=f+g,f,g$. The diagonal cross term is zero. In the beta cross term the translated and reflected copies of $g$ coincide on $\supp f$ at $y=1/2$ by symmetry. For $y=\ell/2$, $\ell\ge2$, both vanish there. Also $g(u)f(u+y)=0$ for every such $y$, and the two reflected cross terms have equal integral. The odd-character projection is the difference of the equal- and opposite-residue indicators on units. For $m$ with $f(m/M)\ne0$ only the two samples $N+m,N-m$ contribute from $g$, both with value one. Thus the real part of the mixed matrix element is $-B_N(b_0)/M$.

Mellin Plancherel for the weighted functional gives
\[
 \|x_v\|^2=\frac1{2\pi}\int_0^\infty
 \frac{(Mu)^{2a_N}}{1+(Mu)^{-2(b-a_N)}}\bigl(v(u)/2-u v'(u)\bigr)^2\dd u .
\]
The weight is bounded by $e^4$ on the supports. For $v=f$, the change of variable $y=Mu$ contributes $M^{-1}$, and the two support constraints give
$H\int(b_0(y)/2-yb_0'(y))^2\dd y\le9P^2$. The corresponding integral for the fixed $g$ is constant. Apply Cauchy--Schwarz to the mixed element and use $M=2N$.
\end{proof}

\section{An odd layer of inclusion--exclusion}\label{sec:sieve}

Fix an odd positive integer $k$. For a positive integer $t$ set
\[
 \begin{gathered}
 z=8t^4,\qquad H=t^{4k},\qquad D=32\cdot16^k,\qquad P=DH,\\
 \mathcal P=\{p:z<p\le2z\},\qquad s=|\mathcal P|,\qquad Q=2\prod_{p\in\mathcal P}p .
 \end{gathered}
\]
Choose smooth functions $u,v$ taking values in $[0,1]$, supported respectively in $[1,4]$ and $[1,2]$, and equal to one on $[2,3]$ and $[7/5,8/5]$, respectively. Set $b_a(x)=u((x-a)/H)v(x/P)$ and
\[
 R(a)=\sum_{P\le m\le2P}\Lambda(m)b_a(m)
 \bigl(\one_{(a-m,Q)=1}-\one_{(a+m,Q)=1}\bigr).
\]
The product rule gives an absolute $V\ge1$ such that $b_a/V$ is a test with interval length $3H$ whenever $20P<N$.

\begin{lemma}\label{lm:row}
For all sufficiently large $t$ there is an even integer $a\in[0,2P]$ with $R(a)\gg_k t^{3k}$.
\end{lemma}
\begin{proof}
Sum over the even rows $a$. Even $m$ contribute zero. For odd $m$ write $h=m-a$, $w(h)=u(h/H)$ and $v_m=\one_{m\ {\rm odd}}\Lambda(m)v(m/P)$. The result is
\[
 \sum_{P\le m\le2P}v_m
 \sum_{\substack{0<h<4H+1\\ h\ {\rm odd}}}w(h)
 \bigl(\one_{(h,Q)=1}-\one_{(2m-h,Q)=1}\bigr).
\]
On the support of $w$ every such $h$ corresponds to an allowed row. Expand both indicators over subsets $T\subseteq\mathcal P$, writing $d_T=\prod_{p\in T}p$. Eventually $4H+1\le z^k$ and $4P<z^{k+1}$. Consequently $k$ distinct band primes cannot divide $h$ together, and $k+1$ cannot divide $2m-h$. In layer $k$ only $d_T\mid2m-h$ survives. Its coefficient is $+1$ because $k$ is odd.

For $|T|<k$, the two counts of odd $h$ are residue classes $d_T$ and $2m-d_T$ modulo $2d_T$. The prefix counts differ by at most two. Abel summation bounds their weighted difference by twice an endpoint value plus twice the discrete total variation of $w$, hence by an absolute constant. There are at most $k(s+1)^{k-1}$ such subsets and $\sum v_m\le(P+1)\log(2P)$.

For a top-layer subset, $d=d_T\le(2z)^k=P/32$. Take an odd multiple $y$ of $d$ within distance $d$ of $3P$. Each prime $m$ in
\[
   (y+2H)/2<m\le(y+3H)/2
\]
contributes $\log m$ to that subset, since $h=2m-y$ is odd, $2H\le h\le3H$, and $7P/5\le m\le8P/5$. This interval has length $P/(2D)$. The prime number theorem $\vartheta(x)\sim x$ gives mass at least $H/4$ uniformly in $d$ for fixed $k$. Indeed both endpoints lie in $[P,2P]$, on which the endpoint errors are uniformly $o(P)$. Thus
\[
 \sum_{a\ {\rm even}}R(a)\ \ge\
 \binom{s}{k}\frac H4-O\bigl(k(s+1)^{k-1}(P+1)\log(2P)\bigr).
\]
The same prime number theorem gives $s\log(2z)\ge z/2$ eventually, so $s\ge t^3$. For fixed $k$, $\binom{s}{k}\gg_k s^k$, $H=P/D$ and $\log(2P)=o(s)$. Divide the resulting lower bound by at most $2P+1$ rows.
\end{proof}

\section{From residue classes to an even modulus}\label{sec:transfer}

Retain the broad weight $g$ of Lemma~\ref{lm:transfer}. For a residue class $v\bmod Q$ and a new sampling scale $M$, write
\[
 A(v)=\sum_{\substack{n<M\\n\equiv v\bmod Q}}\Lambda(n)g(n/M),
 \qquad U=\frac1{\varphi(Q)}\sum_{\substack{n<M\\(n,Q)=1}}\Lambda(n)g(n/M).
\]
For $(v,Q)=1$ and $G(Q)$ one has
\begin{equation}\label{eq:ap}
 |A(v)-U|\ll_g M^{3/4}Q^2 .
\end{equation}
To see this, let $F(u)=\sqrt u\,g(u)$ and $\chi$ be a nonprincipal character modulo $Q$. Mellin inversion for $\sum_n\Lambda(n)\chi(n)n^{-1/2}F(n/M)$ on $\Rea s=1$ has integrand $M^s\widehat F(s)(-L'/L)(1/2+s,\chi)$. On $1/4\le\Rea s\le1$ the logarithmic derivative is holomorphic under $G(Q)$ and has the usual bound $O(\log^2(Q(2+|\Im s|)))$. This is the proved strip estimate for Dirichlet logarithmic derivatives used in the Li--Liu construction \cite{LL,LLLean}. Four integrations by parts give $\widehat F(s)\ll_g(1+|\Im s|)^{-4}$. After a shift to $\Rea s=1/4$, replacing the logarithmic bound by $O(Q^2(1+|\Im s|)^2)$ gives $O_g(M^{1/4}Q^2)$. Multiply by $\sqrt M$ and average the nonprincipal characters by orthogonality to obtain \eqref{eq:ap}. For $4Q\le M$ and sufficiently large $M$, the primes in $(7M/16,9M/16]$ exceed $Q$ and contribute at least $M/16$ by the prime number theorem. Hence $U\ge M/(16\varphi(Q))$.

For $0<m\le2P$, $8P\le M$ and $4Q\le M$ this implies
\begin{equation}\label{eq:translate}
 \sum_{\substack{N<M,\,N\equiv a\bmod Q\\(N,m)=1}}g(N/M)\Lambda(N\pm m)
  =U\one_{(a\pm m,Q)=1}+O_g\bigl(M^{3/4}Q^2+\sqrt M\log M+P\log(2M)\bigr).
\end{equation}
Indeed moving the weight from $N/M$ to $(N\pm m)/M$ costs $O_g(m\log(2M))$ by the mean value theorem. Its support stays away from the endpoints. After translation $(N\pm m,m)=(N,m)$. Every prime in the new support exceeds $m,Q$. A discarded term in a unit class, or a surviving term in a nonunit class, must therefore be a proper prime power. Their total weight is bounded by $\psi(M)-\vartheta(M)\le2\sqrt M\log M$. Apply \eqref{eq:ap} to the remaining unit-class sum.

Set $r=2^{t^9}$ and $M=r^{64}$. Eventually $P,Q\le r$, since $P$ is polynomial in $t$ and $\log_2 Q\le 1+2z(2z+1)\le t^9$. Multiplying the error in \eqref{eq:translate} by $2(P+1)\log(2P)$ gives at most $M/(32Q)$ for large $t$. For example the error for a single sign is at most $5r^{51}$ after absorbing fixed constants and logarithms, and $640r^{54}\le r^{64}$ eventually. Take the even $a$ of Lemma~\ref{lm:row}. Subtract the two signs, multiply by $\Lambda(m)b_a(m)$, and sum over $m$. For $t$ large enough that $R(a)\ge1$ this gives
\[
 \sum_{\substack{N<M\\N\equiv a\bmod Q}}g(N/M) B_N(b_a)
 \ge U R(a)-M/(32Q)\ge M R(a)/(32Q).
\]
There are at most $2M/Q$ candidate centres, and $0\le g\le1$. Some even $N\in(M/4,M)$ therefore satisfies $B_N(b_a)>R(a)/128$. At that centre $20P<N$ eventually. Lemma~\ref{lm:transfer} and the normalized test $b_a/V$ give
\begin{equation}\label{eq:winning}
  N\|T_N\|^2\gg \frac{(3H)B_N(b_a/V)^2}{P^2}
       \gg_k t^{2k},\qquad \log N\le64t^9\log2 .
\end{equation}
For any given $d$ take odd $k$ with $2k>9d$ and let $t$ tend to infinity.

\begin{proof}[Completion of the proof of Theorem~\ref{th:growth}]
One has $1\le\SSg(N)\le(\log N)^2$ eventually. For the lower bound, completing the factors for $p-1$ in $C_2$ to all integers $n\ge2$ gives the smaller product $\prod_{n=2}^m(1-n^{-2})=(m+1)/(2m)$ in the limit. For the upper bound order the odd prime factors as $p_i\ge i+2$. Their correction factors telescope upwards to at most $1+\omega(N)$, where $2^{\omega(N)}\le N$.
Hypothesis R would give $N\|T_N\|^2\le(\log N)^5$, contradicting \eqref{eq:winning}. At width $\eta=1/(4(\log N)^j)$, criterion~\eqref{eq:dict} would imply $N\|T_N\|^2\le(4/3)(\log N)^{j+5}$. Choosing $d=j+6$ gives a strict violation at arbitrarily large $N$ and a real zero in $(1-\eta,1)$. A larger prescribed width also fails there. Choose the $i$th modulus beyond its predecessor with $j=i$. For fixed $d$ the quantity $(\log N_i)^d(1-x_i)$ is eventually at most $1/(4\log N_i)$.
\end{proof}

\section{Arithmetic records and sparse prescriptions}\label{sec:records}

For $n>1$ define
\[
 h_j(n)=\sqrt{\frac{n\SSg(n)}{4(\log n)^{j+3}}},\qquad j=0,1,\ldots
\]
and give $h_j$ value one at $0,1$. Then $\eta_{h_j}(n)=1/(4(\log n)^j)$ for $n>1$ and $h_j(n)\to\infty$. Write $A_d(N)$ if some test of Section~\ref{sec:entry} satisfies $HB_N(b_0)^2>P^2(\log N)^d$. Equation~\eqref{eq:winning}, at the stage before applying the operator bound, proves $A_d$ at arbitrarily large even moduli for each fixed $d$.

Put $r_{-1}=2$. Take $r_j$ to be the least even $N>r_{j-1}$ satisfying $A_{j+8}(N)$ and $h_j(n)\ge j$ for every integer $n\ge N$. Both requirements can be met by the preceding paragraph. Set
\[
 s(n)=\max\bigl(\{0\}\cup\{j:r_j\le n\}\bigr),\qquad h_*(n)=h_{s(n)}(n).
\]
Thus $s(r_j)=j$ and $s(n)\to\infty$. For $n\ge r_0$ one has $h_*(n)\ge s(n)$, and $\eta_{h_*}(n)=1/(4(\log n)^{s(n)})$ eventually. In particular $h_*$ is admissible and $(\log n)^d\eta_{h_*}(n)\to0$ for each fixed $d$.

\begin{lemma}\label{lm:records}
At every sufficiently late record $N=r_j$, the exclusion \eqref{eq:free} for $h_*$ fails.
\end{lemma}
\begin{proof}
If the exclusion holds there, the norm criterion \eqref{eq:dict}, $1-\eta_{h_*}\ge3/4$ and $\SSg(N)\le(\log N)^2$ imply $N\|T_N\|^2\le(\log N)^{j+6}$ once $\log N\ge2$. The threshold here is independent of $j$. But $A_{j+8}(N)$ and \eqref{eq:transfer} give $(\log N)^{j+8}<C N\|T_N\|^2$. Both inequalities cannot hold when $(\log N)^2\ge C$.
\end{proof}

\begin{proof}[Proof of Theorem~\ref{th:sparse}]
Positivity, divergence and the eventual bound $\eta\le1/4$ hold for each of $h_*,h_0$ and pass to the piecewise choice. At $r_{p(j)}$ the two prescriptions $\widetilde h$ and $h_*$ have the same value and width, so Lemma~\ref{lm:records} gives a violation for each sufficiently large $j$. Since $r_{p(j)}\ge p(j)\ge2^j$, a record in $[1,X]$ requires $j\le\lfloor\log_2 X\rfloor$. Finally the logarithmic base has $\eta_{\log}(n)=(\log n)^5/(n\SSg(n))\le(\log n)^5/n$. Multiplying either width in the piecewise choice by $(\log n)^d$ gives a sequence tending to zero.
\end{proof}

\section{Finite kernel tests for Hypothesis S}\label{sec:samples}

For $n=0,1,\ldots$, set $z_n=1/9+i/(n+1)$. These points lie in every strip considered here and converge to the interior point $1/9$. Set $K_{ij}=K_N(z_i,z_j)$ and $G_{ij}=G_{a_N,b}(z_i,z_j)$, where the strip reproducing kernel is
\[
 G_{a,b}(s,z)=\frac{1}{4(b-a)}
     \csc\!\left(\frac{\pi(s+\bar z-2a)}{2(b-a)}\right).
\]
For finite coefficient vectors $c,d$ indexed by the nonnegative integers define
\[
 F_A(c,d)=\sum_{i,j}\overline{d_i}\,c_j A_{ij},
 \qquad g_N(c)=\sqrt{\Rea F_G(c,c)} .
\]
Write $\mathcal M_N(B)$ for the collection of inequalities
$|F_K(c,d)|\le B g_N(c)g_N(d)$ for every such $c,d$. The support size may vary with $N$.

\begin{theorem}\label{th:samples}
For each eligible even $N>e^{10}$ and $B\ge0$,
\[
 \mathcal M_N(B)\quad\Longleftrightarrow\quad \|T_N\|\le B .
\]
Consequently, for each admissible $h$, $\mathrm S(h)$ is equivalent to the eventual validity on eligible even integers of
\begin{equation}\label{eq:sampleS}
 \mathcal M_N\!\left(\frac{\delta_N}{2\sqrt{\eta_h(N)(1-\eta_h(N))}}\right).
\end{equation}
\end{theorem}
\begin{proof}
Take $v_c=\sum_j c_j k_{z_j}$. A Hardy vector orthogonal to every $k_{z_j}$ represents a holomorphic function vanishing at every $z_j$. The identity theorem on the connected strip makes it identically zero. Injectivity of the Fourier--Laplace realization of the Hardy space then gives the zero vector. The span of the evaluation vectors is therefore dense. Moreover
$F_K(c,d)=\langle T_Nv_c,v_d\rangle$ and $g_N(c)=\|v_c\|$ by the two reproducing identities. The operator bound gives $\mathcal M_N(B)$ by Cauchy--Schwarz. Conversely the form bound on the dense span extends to all vectors by continuity. Taking the second vector to be $T_Nv$ gives $\|T_Nv\|\le B\|v\|$. Formula~\eqref{eq:sampleS} follows from \eqref{eq:dict}, including the case of equality.
\end{proof}

\begin{theorem}\label{th:exist}
There exists an admissible $h$ satisfying $\mathrm S(h)$ if and only if
$\|T_N\|/\SSg(N)\to0$ along eligible even integers. Equivalently, every admissible prescription fails if and only if there exist a fixed $\varepsilon>0$ and arbitrarily large eligible even $N$ for which some finite $c,d$ satisfy
\[
             |F_K(c,d)|>\varepsilon\SSg(N)g_N(c)g_N(d).
\]
\end{theorem}
\begin{proof}
At the prescribed width, \eqref{eq:dict} reads
\begin{equation}\label{eq:relative}
   2h(N)\sqrt{1-\eta_h(N)}\,\|T_N\|\le\SSg(N).
\end{equation}
Necessity follows since $h$ diverges and $\eta_h\le1/4$. Conversely put $a_N'=\|T_N\|$ on eligible even $N$, extend by zero elsewhere, and take for $N\ge4$
\[
 h(N)=
 \begin{cases}
 \min\{\log N,\SSg(N)/(2a_N')\},&a_N'>0,\\
 \log N,&a_N'=0 .
 \end{cases}
\]
Extend $h$ positively to smaller indices. The limit hypothesis implies divergence of this single function along all even integers. Also $h\le\log N$ ensures the required bound on the width. By construction $2h(N)a_N'\le\SSg(N)$, which suffices for \eqref{eq:relative}. Negating the limit gives a fixed positive relative level with arbitrarily large strict violations. Theorem~\ref{th:samples} expresses each violation by finite coefficients.
\end{proof}

\section{How large must a shift test be?}\label{sec:limit}

For $h(N)=\log N$ the required upper bound in \eqref{eq:relative} has size $\SSg(N)/\log N$. Normalize a single test by
$D_N=HB_N(b_0)^2/(P^2N)$, so \eqref{eq:transfer} gives $\|T_N\|^2\ge D_N/C$. A lower bound of the required order from this inequality alone must have $D_N(\log N)^2/\SSg(N)^2$ bounded away from zero.

\begin{proposition}\label{pr:cap}
For any even $N\ge4$ and any test, without a GRH assumption,
\[
 |B_N(b_0)|\le28P\log N,\qquad
 D_N(\log N)^2\le784\,\frac{P}{N}(\log N)^4 .
\]
Consequently $D_N(\log N)^2\to0$ uniformly in all tests with $P^2\le N$. More generally it tends to zero along any sequence with $P=o(N/(\log N)^4)$.
\end{proposition}
\begin{proof}
Both $\Lambda(N-m)$ and $\Lambda(N+m)$ lie in $[0,\log(2N)]$. Since $|b_0|\le1$, their difference yields
$|B_N(b_0)|\le\log(2N)\psi(2P)$, where $\psi(x)=\sum_{n\le x}\Lambda(n)$. Chebyshev's inequality $\psi(x)\le(\log4+4)x$, together with $\log(2N)\le2\log N$ and $\log4+4\le7$, gives the first bound. Square it, multiply by $H/(P^2N)$, and use $H\le P$. The uniform assertion follows from $(\log N)^4/\sqrt N\to0$.
\end{proof}

For the moduli of Section~\ref{sec:transfer}, $P$ grows polynomially in $t$ and $N>2^{64t^9}/4$. The proposition applies even though the same tests prove Theorem~\ref{th:growth}. On the upper-bound side the Hilbert--Schmidt estimates give $\|T_N\|\ll(\log N)^2$ under local GRH \cite[Theorem 6.2]{LL}. By \eqref{eq:dict} they guarantee the width $c\SSg(N)/(N\log N)$ for a fixed $c>0$. The ratio of $\eta_{\log}$ to this width is $(\log N)^6/(c\SSg(N)^2)$ and tends to infinity. We have neither a uniform bound at the threshold \eqref{eq:sampleS} for $h=\log$ nor a cofinal strict violation at that threshold. The logarithmic prescription and the existence assertion of Theorem~\ref{th:exist} remain undecided here.

\section{Prime counts on the first physical scale}\label{sec:sharp}

For $0\le t\le1$ let $P_N(t)$ count ordered pairs of primes $p+q=N$ with $p<tN$. Put $L=\log N$, $a=1/L$ and $\beta=1/8$. Define a positive measure on $[0,1]$ by
\[
 \mu_N=\frac{1}{N\SSg(N)}
 \sum_{\substack{n<N\\(n,N)=1}}\Lambda(n)\Lambda(N-n)\,\delta_{n/N},
 \qquad F_N(t)=\mu_N([0,t)).
\]
For a real $f\in C_c^6(0,1)$ set $R_f(u)=f(u)f(1-u)$, $Df(u)=f(u)/2-u f'(u)$ and
\begin{align*}
 W_N(u)&=\frac{(Nu)^{2\beta}}{1+(Nu)^{2(\beta-a)}},&
 E_N(f)&=\frac{1}{2\pi}\int_0^1 W_N(u)|Df(u)|^2\dd u,\\
 Q_N(f)&=\SSg(N)\left(\int_0^1 R_f(u)\dd u-\int R_f\dd\mu_N\right).
\end{align*}
At scale $M=N$, $E_N(f)$ is the squared norm of the test vector, and $Q_N(f)$ is its quadratic form under $T_N$. To see the arithmetic identity, odd-character orthogonality gives the difference of the congruences $m\equiv n$ and $m\equiv-n\bmod N$ on units. For $0<n,m<N$ the first congruence is the diagonal and the second is $m+n=N$. The model on this range consists of the single reflected integral. The Mellin norm formula gives $E_N$ by changing variables $u=e^{-t}/N$ in the Fourier representation of the Hardy norm.

\begin{theorem}\label{th:sharp}
Along every sequence of eligible even moduli $N\to\infty$, the following limits are equivalent:
\[
 \sup_{0\ne f\in C_c^6(0,1;\mathbb R)}
 \frac{|Q_N(f)|}{\SSg(N)E_N(f)}\longrightarrow0,
 \qquad
 \sup_{0\le t\le1}\left|\frac{L^2P_N(t)}{N\SSg(N)}-t\right|
 \longrightarrow0.
\]
The existence of an admissible $h$ satisfying $\mathrm S(h)$ implies both limits.
\end{theorem}

\begin{lemma}\label{lm:totalvar}
Uniformly for sufficiently large even $N$,
\[
 \left\|\mu_N-\frac{L^2}{N\SSg(N)}
       \sum_{\substack{p+q=N\\p,q\ {\rm prime}}}\delta_{p/N}
       \right\|_{\mathrm{TV}}
 \ll \frac1L+\frac{L^2}{\sqrt N}.
\]
Here the norm of an atomic signed measure is the sum of the absolute masses of its distinct atoms.
\end{lemma}
\begin{proof}
Write $A_N(I)$ for the number of ordered representations with first prime in an interval $I$ of $Y$ consecutive integers in $[0,N]$. We first establish
\begin{equation}\label{eq:sharp-sieve}
 L^2 A_N(I)\ll Y\SSg(N)+\sqrt N L^2.
\end{equation}
Take $z=\lfloor N^{1/64}\rfloor$, $R=z^8$ and $P=\prod_{p\le z}p$. At $p$ the forbidden residues $0,N$ have density $\nu_p=(1+\one_{p\nmid N})/p$. For squarefree $d\mid P$, the Chinese remainder theorem gives density $\nu_d=\prod_{p\mid d}\nu_p$ for $d\mid n(N-n)$. The error on any translated interval is at most $2d$, by separating complete periods from the two ends. The Selberg upper sieve therefore bounds the surviving count by $Y/G(R)+2R^4$, where
\[
 G(R)=\sum_{\substack{d\mid P\\d\le R}}\prod_{p\mid d}\frac{\nu_p}{1-\nu_p}.
\]
Indeed the optimal square weights satisfy $|\lambda_d|\le1$ and vanish for $d>R$. Summing the errors $2[d,e]$ over $d,e\le R$ gives $2R^4$. Normalize the unrestricted summands of $G$ to a probability measure. Each prime is included independently with probability $\nu_p$, so the expectation of $\log d$ is at most $4\log z$: the divisibility of $z!$ by prime powers gives $\sum_{p\le z}\log p/p\le2\log z$. Markov's inequality gives $G(R)\ge\frac12\prod_{p\le z}(1-\nu_p)^{-1}$. Multiplying the local factors and applying Mertens' product estimate gives $\prod_{p\le z}(1-\nu_p)\ll\SSg(N)/\log^2z$. There are at most $2(z+1)$ additional pairs with a summand at most $z$. This proves \eqref{eq:sharp-sieve}. Including proper prime powers, the corresponding sum of $\Lambda(n)\Lambda(N-n)$ over $I$ is $\ll Y\SSg(N)+\sqrt N L^2$, since $\psi(N)-\vartheta(N)\ll\sqrt N$ and each complementary weight is at most $L$.

Let $b_n$ be the indicator that $n$ and $N-n$ are both prime. Telescoping $\log N-\log n$ and interchanging two finite sums gives
\[
 \sum_{n<N} b_n(L-\log n)
 =\sum_{1\le v<N}\log(1+1/v)\sum_{n\le v}b_n
 \ll \frac{N\SSg(N)}{L^2}+\sqrt N L.
\]
For the main term use $(v+1)\log(1+1/v)\le2$ in \eqref{eq:sharp-sieve}. The remaining logarithmic increments sum to $L$. Reflection gives the same bound at the other summand. On a prime pair,
$0\le L^2-\log p\log q\le L[(L-\log p)+(L-\log q)]$.
Removing proper prime powers costs $O(\sqrt N L)$ in unnormalized mass. A prime pair excluded by coprimality must have $p=q=N/2$ and costs at most $L^2$. Division by $N\SSg(N)$ proves the estimate, using $\SSg(N)>1$.
\end{proof}

\begin{proof}[Proof of Theorem~\ref{th:sharp}]
Assume the quadratic limit. For fixed $f$, $W_N(u)\le(Nu)^{2/L}\le e^2$ gives a uniform bound for $E_N(f)$. Thus $\int R_f\dd\mu_N\to\int_0^1R_f\dd u$. Given an arbitrary real $\varphi\in C_c^6(0,1)$, choose a compact smooth $g$ equal to $1$ on the reflection of $\supp\varphi$. Then
\[
 R_{g+\varphi}(u)-R_{g-\varphi}(u)=2\{\varphi(u)+\varphi(1-u)\}.
\]
Both measures are invariant under $u\mapsto1-u$, so their integrals of $\varphi$ converge. The weighted version of \eqref{eq:sharp-sieve} gives
\[
 0\le F_N(t)-F_N(s)\ll t-s+1/N+L^2/\sqrt N
 \qquad(0\le s\le t\le1).
\]
Approximate $[0,t)$ from within by a smooth function with transition intervals of length $2d$ at either end. The displayed bound controls the lost mass by $O(d+1/N+L^2/\sqrt N)$, even at $t=1$. First send $N$ to infinity and then $d$ to zero. Monotonicity and interpolation on a finite mesh give $\sup_t|F_N(t)-t|\to0$. Lemma~\ref{lm:totalvar} transfers the limit to $P_N$.

Conversely, the prime-count limit and Lemma~\ref{lm:totalvar} give $d_N=\sup_t|F_N(t)-t|\to0$. Put $H(t)=\SSg(N)(F_N(t)-t)$, $v(t)=1-e^{-Nt}$ and $B(t)=H(t)-tH(1)$. Integration by parts gives $Q_N(f)=\int_0^1 H R_f'\dd t$. We give the uniform estimates needed when the support of $f$ approaches an endpoint. Since $v(t)\ge Nt/(1+Nt)$ and the first and last intervals of length $1/N$ have no positive atom in their interiors,
\[
 |B(t)|\ll \SSg(N)(d_N+1/N)v(t)v(1-t).
\]
On $[1/N,1-1/N]$ both factors are bounded below. On the two end intervals $B(t)$ is linear and vanishes at the adjacent endpoint, which gives the stated estimate there.

Write $g=vf$, $A=vtf'$ and $J=\int_0^1(vDf)^2\dd t$. Integration by parts, using $v'\ge0$, gives $\int A^2+\tfrac34\int g^2\le J$. Since $0\le tv'\le v$, this also bounds $\|tg'\|_2$ by a constant times $\sqrt J$. Hardy's inequality at the right endpoint applied to $tg$ yields
$\|g/(1-t)\|_2\le\|g\|_2+2\|(tg)'\|_2\ll\sqrt J$.
Cauchy--Schwarz after reflection therefore gives
\[
 \int_0^1 v(t)v(1-t)|R_f'(t)|\dd t\ll J,\qquad
 \int_0^1|R_f(t)|\dd t\ll J.
\]
For the second estimate use $v(t)\ge t/2$ and
$1/(t(1-t))=1/t+1/(1-t)$. Splitting at $Nt=1$ in the energy weight gives $v(t)^2\le2W_N(t)$, hence $J\le4\pi E_N(f)$. Finally $\int_0^1 tR_f'\dd t=-\int_0^1 R_f\dd t$. The two parts of $H=B+tH(1)$ give
\[
 |Q_N(f)|\ll\SSg(N)(d_N+1/N)E_N(f)
\]
uniformly in $f$. This proves the converse. The last implication follows from Theorem~\ref{th:exist} and $|Q_N(f)|\le\|T_N\|E_N(f)$. The other physical scales remain in the existence criterion for the full operator.
\end{proof}

\section{Two-prime counts in large residue classes}\label{sec:congruences}

Write $z=\lfloor\log N/64\rfloor$, $Q=\prod_{p\le z}p$ and
\[
 k_N=\left\lceil\frac{\exp(8\log N\log\log N)}{N}\right\rceil.
\]
For $j\ge0$ set $q=2^jQ$. Let $f,g\in C_c^6(0,1)$ be real.
Define $\mathcal P_q^+(H)$ by summing
$\Lambda(n)\Lambda(m)f(n/M)g(m/M)$ over $0\le n,m<M$ with
$(nm,q)=1$ and $n+m\equiv H\pmod q$. For $\mathcal P_q^-(H)$ impose
$n-m\equiv H\pmod q$. Write $I_f=\int_0^1f(u)\dd u$.

\begin{theorem}\label{th:congruences}
Uniformly for $M\ge N$, $0<H\le 2k_NN$ even, and all $j\ge0$ such
that $q(\log M)^3\le M$ and $G(q)$ holds, one has
\[
  \frac{q\,\mathcal P_q^\pm(H)}{M^2\SSg(H)}=I_f I_g+o_{f,g}(1)
             \qquad(N\to\infty).
\]
For every $M\ge N$ and sufficiently large $N$, there exists a lift
$M/(2\log^3 M)<q\le M/\log^3 M$.
More generally the limit holds along any sequence with $q\le M$ and
$q(6+\log q)^2/M\to0$.
\end{theorem}

\begin{proof}
Let $E_q(H)$ consist of the residues $t\bmod q$ for which
$t(t-H)$ is a unit. Put $e=\#E_q(H)$, $\phi=\varphi(q)$ and
$T=qe/\phi^2$. For $z\ge2$, the prime divisors of $q$ are precisely
those of $Q$. Passing from $Q$ to $q$ multiplies both $e$ and
$\varphi(Q)$ by $2^j$. The Chinese remainder theorem at $Q$ gives
\[
 T=2\prod_{2<p\le z}\left(1-\frac1{(p-1)^2}\right)
       \prod_{\substack{2<p\le z\\p\mid H}}\frac{p-1}{p-2}.
\]
The ratio $T/\SSg(H)$ tends uniformly to one in the stated range.
Indeed $\log H\ll z\log z$. For each fixed $\eta>0$, Mertens' theorem
bounds the reciprocal mass of primes $z<p\le z^{1+\eta}$ by
$\log(1+\eta)+o(1)$. The larger prime divisors of $H$ have reciprocal
sum at most $\log H/(z^{1+\eta}\log z)\ll z^{-\eta}$. Letting
$\eta\downarrow0$ handles the omitted factors at primes dividing $H$.
The other Euler product converges absolutely. In particular $T\ge1/2$
for all sufficiently large $N$.

For a unit residue $t$ let $A_f(t)=\sum_{n\equiv t\bmod q}
\Lambda(n)f(n/M)$, where $0\le n<M$. Write $F=\sum_t A_f(t)$,
over the units, and $D_f(t)=A_f(t)-F/\phi$. Character orthogonality gives
\[
  \phi\sum_{t\in(\mathbb Z/q)^\times}|D_f(t)|^2
    =\sum_{\chi\ne\chi_0}
       \left|\sum_{n<M}\Lambda(n)\chi(n)f(n/M)\right|^2 .
\]
We prove that under $G(q)$ each nonprincipal twist is bounded by
$B_f=C_f\sqrt M(6+\log q)$ for sufficiently large $M$.
For the primitive character inducing $\chi$ put
\[
 Z_\sigma(t)=\sum_\gamma\frac{\sigma}{\sigma^2+(t-\gamma)^2},
 \qquad P_\sigma(t)=\frac{\sigma}{\pi(\sigma^2+t^2)}.
\]
Here $\gamma$ runs through the ordinates of the nontrivial zeros with
multiplicity. Positivity and $P_1*P_\sigma=P_{1+\sigma}$ give
$\int P_1(t)Z_\sigma(t)\dd t=Z_{1+\sigma}(0)\le3Z_{3/2}(0)$
for $0<\sigma\le1$. The zero sum at $3/2$ satisfies
$Z_{3/2}(t)\ll\log(2q(2+|t|))$ by the zero count on unit intervals.
Write $D(t)=(L'/L)(1/2+\sigma+it,\chi)$. Take the difference of the
completed logarithmic derivatives at $1/2+\sigma+it$ and $2+it$ in
the Hadamard product. Under $G(q)$ the real part of the zero sum is
$Z_\sigma(t)-Z_{3/2}(t)$. The Dirichlet series at $2+it$ is bounded.
The difference of the digamma terms is bounded uniformly in $t$ by
its resolvent series, since both gamma arguments have real part at
least $1/4$ and their separation is at most $3/4$. Each omitted
Euler factor costs $O(\log p)$ at these two points. Hence
\[
 \int_{\mathbb R}\frac{|\Rea D(t)|}{(1+|t|)^4}\dd t\ll 6+\log q.
\]

Set $\sigma=1/\log M$ and
$W(t)=M^{\sigma+it}\widehat{\sqrt u\,f(u)}(\sigma+it)$.
Four integrations by parts give $|W(t)|\ll_f(1+|t|)^{-4}$. Mellin
inversion expresses the twist as $-\sqrt M(2\pi)^{-1}\int W D$.
We claim $\int\overline W D=0$. Indeed let
$g(v)=\sqrt{v/M}\,f(v/M)$ and take $g(v)=0$ for $v\le0$.
On this line the Mellin transform of $v^{-2\sigma}g(1/v)$ equals
$\overline W$. For large $M$ its support is a compact subset of
$(0,1)$, so every positive integer sample in the inversion formula
vanishes. The contour move is valid for both compact weights by the
zero-free half-plane and rapid Mellin decay. Conjugating the zero
integral gives $\int W\overline D=0$. Thus $\int W D=2\int W\Rea D$,
which proves the bound for the twist.
Parseval now gives $\sum_t |D_f(t)|^2\le B_f^2$.

Put $U=F/M$ and define $V$ from $g$ in the same way. The map
$\tau(t)=H-t$ for the sum, or $\tau(t)=t-H$ for the difference,
takes $E_q(H)$ injectively into the units. Expanding
$\sum_{t\in E_q(H)}A_f(t)A_g(\tau(t))$ about the two means, then
applying Cauchy--Schwarz on $E_q(H)$, gives
\[
 \left|\frac{q\,\mathcal P_q^\pm(H)}{M^2T}-UV\right|
 \le b_f b_g+|U|b_g+|V|b_f,\qquad
 b_f=\frac{\phi B_f}{M\sqrt e}.
\]
As $\phi^2/e=q/T\le2q$, we have $b_f\ll_f(6+\log q)\sqrt{q/M}$.
This tends to zero under the stated general condition, and is
$O_f((\log M)^{-1/2})$ at the displayed dyadic lift.

Finally $U=I_f+o_f(1)$ uniformly. Partial summation against
$\psi(x)=x+o(x)$ gives the claim before the restriction to units.
Only primes dividing $Q$ and their powers can be removed.
Here $Q\le\sqrt M$, so for large $M$ the primes themselves lie below
the support of $f(n/M)$. The total weight at proper prime powers is
$O(\sqrt M)$. The same argument applies to $V$. Combining these
estimates proves the limit. The bound $Q\le\sqrt M$ also gives
$Q(\log M)^3\le M$ eventually, so successive doublings cross
$M/\log^3M$ and give the stated lift.
\end{proof}

At this lift the support meets $O(\log^3 M)$ possible lines in a
residue class. An exact identity records what remains to isolate any
one of them. For either sign, fix the two weights including the unit
conditions at $q$. Write $P_u(d)$ for their congruence sum at $d\bmod u$
and $G(d)$ for their sum on $n\pm m=d$. For $0\le H\le 2M$ and
$2^rq>3M$, one has
\[
 2^r G(H)=P_q(H)+\sum_{j=0}^{r-1}2^j
 \bigl(P_{2u_j}(H)-P_{2u_j}(H+u_j)\bigr),\qquad u_j=2^jq .
\]
Indeed the two displayed classes at $2u_j$ partition the class at
$u_j$. At the final modulus, a multiple of $2^rq$ equal to
$n\pm m-H$ must be zero. Telescoping gives the identity.

Removing the unit conditions from this single line changes $G(H)$ by
at most
\[
 2\|f\|_\infty\|g\|_\infty(\log M)\lfloor\log_2 M\rfloor\log q .
\]
To see this, every excluded prime power $p^j<M$ divides
$q^{\lfloor\log_2 M\rfloor}$. The identity
$\sum_{n\mid b}\Lambda(n)=\log b$ bounds the total excluded mass on
either coordinate by $\lfloor\log_2 M\rfloor\log q$. Each coordinate
has at most one partner on the line, whose Mangoldt weight is at most
$\log M$. The remaining analytic question concerns cancellation in
the weighted differences of the lifted congruences.

\section{Fourier inversion on a specified line}\label{sec:frequency}

Fix $\sigma\in\{1,-1\}$ and write $d_\sigma(n,m)=n+\sigma m$.
Let $N,M,H,q$ vary along a sequence satisfying the hypotheses of
Theorem~\ref{th:congruences}, with $N\to\infty$,
$q(6+\log q)^2/M\to0$ and $H/M\to x$. Define
\[
 J_\sigma(x;f,g)=\int_0^1 f(u)g\bigl(\sigma(x-u)\bigr)\,\dd u .
\]
Write $G_\sigma(H)$ for the sum of
$\Lambda(n)\Lambda(m)f(n/M)g(m/M)$ over $0\le n,m<M$ with
$d_\sigma(n,m)=H$, including all prime powers. To resolve a single
line inside its congruence class, put $A=\pi M/q$ and
\[
 P(t)=\frac{q}{M^2\SSg(H)}
 \sum_{\substack{0\le n,m<M,\ (nm,q)=1\\d_\sigma(n,m)\equiv H\bmod q}}
 \Lambda(n)\Lambda(m) f(n/M)g(m/M)
 \cos\!\left(t\frac{d_\sigma(n,m)-H}{M}\right).
\]
For $B>0$ let $E(B)=\int_{-A}^{-B}P(t)\,\dd t+
\int_B^A P(t)\,\dd t$. All integrals here are oriented. For each fixed
$B$, the inequality $B<A$ holds eventually.

\begin{theorem}\label{th:frequency}
Along every such sequence,
\[
 \frac{G_\sigma(H)}{M\SSg(H)}\longrightarrow J_\sigma(x;f,g)
 \quad\Longleftrightarrow\quad
 \lim_{B\to\infty}\limsup_{N\to\infty}|E(B)|=0 .
\]
When $M=H$ and $G(M)$ holds, the tail condition for $\sigma=1$ and
every pair of real tests $f,g\in C_c^6(0,1)$ is equivalent to
both first-scale limits of Theorem~\ref{th:sharp} at the moduli $M$.
In particular the original existence assertion implies this family
of signed tail estimates.
\end{theorem}

\begin{proof}
For an integer $k$, integration of $\cos(tkq/M)$ from $-A$ to $A$
gives $2A$ when $k=0$ and zero otherwise. Thus the integral of $P$
over this interval equals $2\pi G_\sigma(H)/(M\SSg(H))+o(1)$.
The error is bounded by the unit-removal estimate at the end of the
preceding section, divided by $M\SSg(H)$ and multiplied by $2\pi$.

Write $\widehat f(t)=\int_\mathbb R f(u)e^{-itu}\,\dd u$, extending
the tests by zero. At each fixed $t$, apply
Theorem~\ref{th:congruences} to the products of $f$ and $g$ with
$\cos(tu)$ and $\sin(tu)$. Expanding the cosine of the difference
shows
\[
 P(t)\longrightarrow
 F(t)=\Rea\!\left(e^{itx}\widehat f(t)\widehat g(\sigma t)\right).
\]
For each test choose a nonnegative compact smooth majorant for its
absolute value. The congruence estimate on the two majorants gives
a common bound for $|P(t)|$ at every $t$ and all sufficiently large
indices. Dominated convergence applies on every fixed interval
$[-B,B]$. Two integrations by parts give
$\widehat g(t)=O_g((1+|t|)^{-2})$, so $F$ is integrable.
Fourier inversion for $g$ and Fubini's theorem now give
\[
 (2\pi)^{-1}\int_\mathbb R F(t)\,\dd t=J_\sigma(x;f,g).
\]
For each fixed $B$ the difference between the full integral and the
central integral is exactly $E(B)$. The triangle inequality, first
along the arithmetic sequence and then as $B\to\infty$, proves both
implications.

For $M=H$ the archimedean integral uses $g(1-u)$. Imposing
coprimality with $M$ on the exact line costs $O_{f,g}(\log^3 M)$.
The resulting normalized count is
$\int f(u)g(1-u)\,\dd\mu_M(u)$ with the measure of
Section~\ref{sec:sharp}. Every smooth interior test $f$ has the form
$f(u)g(1-u)$ for a smooth $g$ equal to one on the reflection of its
support. Thus convergence on these products is precisely convergence
on all smooth interior tests. The endpoint and energy estimates in
the proof of Theorem~\ref{th:sharp} give the asserted equivalences.
\end{proof}

\section*{Formal verification}
The formal proofs use the residual operator and the hypotheses of the Li--Liu formalization \cite{LLLean}. Lean, lean4lean, con-leche and nanoda independently accepted the proof export for the spectral obstruction and the common export for Theorems~\ref{th:sharp}, \ref{th:congruences} and~\ref{th:frequency}. The latter lean4lean replay checked 84,362 declarations. The exports contain three axiom declarations, for propositional extensionality, classical choice and quotient soundness \cite{formal}.

A separate verified-mode replay of the Dirichlet GRH certificate \cite{grh} accepted 1,468,418 declarations and all five challenge statements. The primitive-character proposition and seven definitions, including the $L$-function and primitivity, agree across the two formal libraries \cite{formal}.

\begingroup\small
\begin{thebibliography}{9}
\bibitem{LL} J. Li and J. Liu, \emph{A conditional proof of the Goldbach conjecture}, 13 November 2026. \href{https://mrc.sdu.edu.cn/english/info/1018/1514.htm}{Shandong University, Research Center for Mathematics}.
\bibitem{YL} Y. Liu, \emph{Failure of the Li--Liu spectral hypothesis R under Dirichlet GRH}, 21 November 2026. \href{https://hexagonmath.org/2611.00788v1}{Hexagon 2611.00788v1}.
\bibitem{LLLean} \emph{Conditional Goldbach in Lean}, revision \texttt{c55db645}. \href{https://github.com/subfish-zhou/conditional-goldbach-lean/tree/c55db645f3b4aed00374fe4aa1fcee2428c88c4b}{Formal definitions and proofs}.
\bibitem{grh} L. Alp\"oge, \href{https://alpo.ge/remark.pdf}{Announcement of a proof of the Riemann hypothesis}, 20 November 2026. A. V. Sutherland, \emph{Verification files for the Lean proof of the Riemann hypothesis}, updated 25 November 2026. \href{https://math.mit.edu/~drew/grh/}{Certificates and challenge statements}.
\bibitem{formal} Codex, \emph{The Li--Liu spectral hypotheses in Lean}, 2026. \href{https://github.com/RaymondVoldemortFu/li-liu-spectral-hypotheses-lean/tree/046b4e9}{Proof sources, revision \texttt{046b4e9}}. \href{https://github.com/RaymondVoldemortFu/li-liu-spectral-hypotheses-lean/tree/main/verification}{Verification records}.
\end{thebibliography}
\endgroup
\end{document}
